\section{Architecture}
\label{sec:architecture}

The CPU backend is history-based and rayon-parallel: each worker
thread pulls a fresh particle from a shared source bank, walks it
to death, accumulates per-batch tallies in a worker-local sink, and
publishes the sinks back to the eigenvalue driver via a
\texttt{fold().reduce()} reduction. The GPU backend is event-batched:
particles live as Structures-of-Arrays in device memory, advance
through a fixed pipeline of kernels (\texttt{find\_cell\_batch},
\texttt{trace\_step\_batch}, \texttt{multi\_step\_walk}, then the
collision-sampling stages), and write tallies via
\texttt{atomicAdd(double*, double)} on Ampere-or-newer hardware.

Both paths consume the same data structures upstream of the
runner: the recursive \texttt{Geometry} (universes, hexagonal and
rectangular lattices, BVH per universe), the \texttt{XsProvider}
trait (one of \texttt{Table}, \texttt{Svd}, \texttt{HybridSvdWmp},
\texttt{HybridTableWmp}), the \texttt{ThermalScatteringData}
registry, and the per-material atom-density tables. They diverge
at the \texttt{EigenvalueRunner} trait, which has two
implementations: \texttt{CpuRunner} (rayon) and \texttt{CudaRunner}
(cudarc + NVRTC).

\paragraph{The shared trait.} \texttt{EigenvalueRunner} exposes a
single method \texttt{run(\&self, config: SimConfig) ->
OutcomeResult}. Every benchmark, every test, and every Python
entry point routes through this method; the choice of backend is a
runtime variable, and the scene JSON corpus does not branch on it.
This is how we run the same 7-case heu-comp-inter-003 family
through CPU and GPU back-to-back with one CLI flag flip, and how
\texttt{tests/cuda\_runs.rs} runs the ICSBEP regression substrate
under \texttt{-\,-features cuda} without duplicating any scene
setup code.

\paragraph{Where the GPU diverges in physics.} A handful of
features are not yet on the device:

\begin{itemize}
    \item GPU survival biasing and Russian roulette. The CPU path
        runs both ($w_\text{min}=0.25$, $w_\text{survive}=1.0$);
        the measured FOM gain on PWR is $4.5\,\times$. The GPU
        runs analog. This is variance-only, not a $k$-eff bias.
    \item GPU discrete S($\alpha$,$\beta$) inelastic (NJOY iwt
        modes 0 and 1). The OpenMC HDF5 distribution emits every
        TSL as \texttt{incoherent\_inelastic} (continuous),
        which is exactly what the device sampler consumes; the
        upload path \texttt{panic!}s if the host hands it
        \texttt{InelasticDist::Discrete} so silent breakage is
        impossible.
    \item Per-cell \texttt{Mat3} rotation. \texttt{GpuRecursiveContext::build}
        errors out loudly if any cell sets
        \texttt{rotation = Some(...)}. No ICSBEP scene currently
        uses it, but the rejection is explicit rather than
        silent.
\end{itemize}

\paragraph{Process-wide shared context.}
\texttt{GpuTransportContext::shared()} returns an
\texttt{Arc<GpuTransportContext>} from a process-wide
\texttt{OnceLock}. First call pays the
NVRTC compile $+$ CUDA init cost ($\sim$30--150~ms on
RTX~A1000); every subsequent caller --- across cases in an ICSBEP
sweep, across PyO3 entry points, across rayon worker threads ---
gets the same Arc. This is the foundation the device-buffer
caches sit on; a fresh context per case (the prior pattern) had
empty caches on every entry. The shared context is also what
makes the \texttt{Arc::as\_ptr}-keyed caches work across cases:
when \texttt{material\_resolve::thermal\_cache} on the host side
returns the same \texttt{Arc<ThermalScatteringData>} to two
different cases, the GPU side gets the same pointer and the cache
hits.
